\documentclass[11pt]{article}

\usepackage[margin=1in]{geometry}
\usepackage{amsmath}
\usepackage{graphicx}
\usepackage{hyperref}
\usepackage{xcolor}
\usepackage{booktabs}
\usepackage{enumitem}
\usepackage{tcolorbox}
\usepackage{parskip}

\hypersetup{
  colorlinks=true,
  linkcolor=blue!60!black,
  urlcolor=blue!60!black,
  citecolor=blue!60!black,
}

\newcommand{\degree}{$^\circ$}

\definecolor{exposure}{HTML}{C0392B}
\definecolor{sensitivity}{HTML}{D35400}
\definecolor{adaptive}{HTML}{16A085}

\newtcolorbox{tipbox}{
  colback=blue!5, colframe=blue!40!black,
  title=\textbf{Tip}, fonttitle=\color{white},
  boxrule=0.5pt, arc=2pt
}
\newtcolorbox{cautionbox}{
  colback=orange!5, colframe=orange!60!black,
  title=\textbf{Worth Knowing}, fonttitle=\color{white},
  boxrule=0.5pt, arc=2pt
}

\title{San Diego Heat Risk Explorer\\
       \large A User's Guide to the Interactive Risk Mapping Tool}
\author{}
\date{\today}

\begin{document}
\maketitle

\tableofcontents
\vspace{1em}
\hrule
\vspace{1em}

%% ==================================================================
\section{What Is This Tool?}
%% ==================================================================

The \textbf{San Diego Heat Risk Explorer} is an interactive map that lets you
build your own heat risk index for San Diego County. You decide which factors
matter --- temperature, air conditioning access, chronic disease rates, tree
canopy, poverty, and more --- and the tool combines them into a neighborhood-level
risk map you can explore, compare, and share.

It is designed for:
\begin{itemize}
  \item \textbf{Public health staff} trying to identify communities most at
        risk during a heat wave.
  \item \textbf{City planners and sustainability teams} prioritizing where to
        plant trees, install cooling centers, or coat roofs.
  \item \textbf{Community advocates} making the case for investment in
        specific neighborhoods.
  \item \textbf{Students, journalists, and curious residents} who want to
        understand how heat risk varies across the county.
\end{itemize}

You do \textbf{not} need to know statistics or GIS to use it. Everything runs
in your web browser --- no software to install, no data to download.

\subsection{Quick Start (60 seconds)}

\begin{enumerate}[leftmargin=*]
  \item Open \texttt{figures/main/diy\_pca.html} in any modern web browser.
  \item Leave the default variable selections as they are.
  \item Click the blue \textbf{Compute PCA} button.
  \item Look at the map --- darker red means higher overall heat risk.
  \item Click any census tract to see which factors drove its risk score.
\end{enumerate}

That's it. The rest of this guide explains what you are looking at, how to
change the inputs, and how to interpret the results responsibly.

%% ==================================================================
\section{The Three Pillars of Heat Risk}
\label{sec:pillars}
%% ==================================================================

Heat doesn't harm everyone equally. A 100\degree\,day hits a shaded,
air-conditioned apartment in Carlsbad very differently from a second-floor
unit in El Cajon with no AC and a 90-year-old resident inside. To capture
that reality, the tool organizes information into three pillars.

\subsection*{\textcolor{exposure}{Pillar 1 --- Exposure}}
\textit{How hot does it actually get here?}

Exposure measures the physical heat load a place experiences. The tool
includes nine exposure variables drawn from weather stations across the
county:

\begin{itemize}[leftmargin=*]
  \item \textbf{Heat:} average daily high, 90th and 95th percentile highs,
        annual maximum temperature, fraction of days above extreme thresholds.
  \item \textbf{Nighttime:} average overnight low, 90th percentile overnight
        low, fraction of nights with no cooling relief.
  \item \textbf{Humidity \& Wind:} average relative humidity.
\end{itemize}

\begin{tipbox}
Nighttime temperatures matter as much as daytime peaks. The body needs
cool nights to recover; persistent warm nights are a strong predictor of
heat-related illness, especially in older adults.
\end{tipbox}

\subsection*{\textcolor{sensitivity}{Pillar 2 --- Sensitivity}}
\textit{Who lives here, and how vulnerable are they?}

Sensitivity captures population characteristics that make people more likely
to be harmed when temperatures rise. It draws on three data sources:

\begin{itemize}[leftmargin=*]
  \item \textbf{Chronic disease prevalence} from the CDC PLACES program ---
        high blood pressure, asthma, heart disease, COPD, diabetes, stroke,
        cognitive difficulty, depression.
  \item \textbf{Demographics and social vulnerability} from the CDC Social
        Vulnerability Index (SVI) --- age 65+, children under 18, disability,
        limited English, poverty, unemployment, housing cost burden, lack of
        a vehicle, and more.
  \item \textbf{Environmental health} --- life expectancy, air pollution,
        ozone, fine particulate matter, diesel exhaust.
\end{itemize}

\subsection*{\textcolor{adaptive}{Pillar 3 --- Adaptive Capacity}}
\textit{What resources does this community have to cope?}

Adaptive capacity is the counterweight to sensitivity: the ability of a
neighborhood to protect itself from heat. Variables include:

\begin{itemize}[leftmargin=*]
  \item \textbf{Air conditioning access} --- probability a home has AC,
        probability of no AC, probability of central AC.
  \item \textbf{Built environment} --- tree canopy coverage, park access,
        share of pre-1980 housing (generally less insulated), renter share,
        green space percentage.
  \item \textbf{Socioeconomic indicators} from the Healthy Places Index ---
        income, education, healthcare access, transportation, housing
        quality, health insurance coverage, homeownership, social cohesion,
        voter participation, access to supermarkets.
\end{itemize}

\begin{cautionbox}
Adaptive capacity is \textit{protective}: more of it means lower risk. The
tool automatically flips the sign when combining pillars, so you don't have
to think about it.
\end{cautionbox}

%% ==================================================================
\section{How the Tool Combines Everything}
\label{sec:pca}
%% ==================================================================

Once you've selected which variables to include, the tool has to combine
them into a single risk score per census tract. The technique it uses is
called \textbf{Principal Component Analysis}, or PCA. You don't need to
understand the math --- but a plain-English version helps you trust the output.

\subsection{PCA in Plain English}

Imagine you have 15 different measurements of every neighborhood: temperature,
poverty rate, asthma rate, AC access, and so on. Many of these measurements
are correlated --- poor neighborhoods tend to have older housing, less tree
canopy, and higher chronic disease rates all at once. Treating them as 15
independent numbers would \textit{double-count} those overlapping patterns.

PCA's job is to find the hidden patterns that explain the most variation
across neighborhoods. It says, in effect:
\begin{quote}
``Most of what distinguishes these tracts from each other can be captured
by just a few underlying patterns. Here they are, ranked by importance.''
\end{quote}

The first pattern (called \textbf{PC1}) captures the biggest source of
variation. The second (PC2) captures the next biggest, and so on. The tool
keeps only the patterns that matter (an eigenvalue threshold called the
\textit{Kaiser criterion}) and weights them by how much variation they explain.

\subsection{From Patterns to Risk Scores}

For each pillar, the tool:
\begin{enumerate}[leftmargin=*]
  \item Puts every variable on the same scale (z-scores).
  \item Finds the main patterns with PCA.
  \item Keeps the important ones.
  \item Combines them into a single pillar score, weighted by importance.
  \item Ranks every tract on a 0--100 scale within that pillar.
\end{enumerate}

Then the three pillar scores are combined into a final composite:
\[
\text{Composite Risk} = \frac{\text{Exposure} + \text{Sensitivity} + (100 - \text{Adaptive Capacity})}{3}
\]

Subtracting adaptive capacity from 100 flips it the right direction ---
more resources means less risk. The result is a single number from 0
(lowest risk) to 100 (highest risk) for every census tract.

%% ==================================================================
\section{Using the Interface}
\label{sec:interface}
%% ==================================================================

The viewer has three regions: the \textbf{left panel} (inputs), the
\textbf{map} (outputs), and the \textbf{results summary} (what just
happened).

\subsection{The Left Panel: Choosing Your Inputs}

\subsubsection*{Climate Scenario}
Four buttons at the top let you pick a time period:
\begin{itemize}[leftmargin=*]
  \item \textbf{Present Day} --- observed weather station data.
  \item \textbf{SSP2-4.5} --- moderate future warming, 2035--2064 (a
        plausible middle-of-the-road emissions path).
  \item \textbf{SSP3-7.0} --- high warming.
  \item \textbf{SSP5-8.5} --- very high warming (sometimes called
        ``business as usual'' fossil-fuel-intensive).
\end{itemize}

Choosing a future scenario shifts the temperature variables upward using
downscaled climate model output (LOCA2 ensemble mean). Sensitivity and
adaptive capacity are held at their present-day values --- we don't
pretend to know who will live in each tract in 2050.

\begin{cautionbox}
Not every temperature variable has a future projection available. Those
without one are shown with a gray ``present-day'' badge and simply stay
at today's value under future scenarios.
\end{cautionbox}

\subsubsection*{Variable Checkboxes}
Below the scenario buttons, each pillar has its own color-coded section
(\textcolor{exposure}{red}, \textcolor{sensitivity}{orange},
\textcolor{adaptive}{teal}) with subsection headings like ``Heat,''
``Nighttime,'' ``Demographics,'' etc.

Next to each variable is a checkbox. \textbf{Checked means the variable is
included} in the PCA. Uncheck anything you don't want to use. A small
\texttt{src} link opens the original data source in a new tab --- use
this to understand exactly where a number comes from.

\subsubsection*{Compute PCA Button}
The large blue \textbf{Compute PCA} button re-runs the analysis with
your current selections. Nothing updates automatically --- you have to
click it. This is deliberate: it lets you experiment freely without the
map thrashing every time you toggle a checkbox.

\subsection{The Map}

The map shows census tracts colored by the currently selected layer.
Darker colors always mean higher values on the chosen metric.

\subsubsection*{Layer Buttons}
Below the left panel, you'll find a list of layers you can switch between:
\begin{itemize}[leftmargin=*]
  \item \textbf{Composite Risk} --- the overall 0--100 risk score.
  \item \textbf{Dominant Driver} --- for each tract, which pillar is
        pushing its risk the highest. Useful for answering ``why is this
        neighborhood high risk?''
  \item \textbf{Individual Pillar Principal Components} --- e.g.,
        ``Exposure PC1 (67\%)'' means the biggest exposure pattern, which
        explains 67\% of the variance across tracts. Selecting one of
        these shows the spatial pattern for that component alone.
\end{itemize}

\subsubsection*{Overlays}
Checkboxes below the layer buttons outline specific city jurisdictions ---
San Diego, Chula Vista, La Mesa, and ``other cities.'' Use these to focus
on a particular jurisdiction's boundary.

\subsubsection*{Tract Details}
Click any tract on the map to see a popup with its numeric scores for
each pillar, its composite risk, and its dominant driver.

\subsection{The Results Summary}

After you click \textbf{Compute PCA}, a small results box appears
listing, for each pillar, how many principal components were retained
and how much of the total variance they explain. For example:
\begin{quote}
\textbf{Exposure}: 1 PC retained (87\% variance explained)\\
\textbf{Sensitivity}: 2 PCs retained (74\% variance explained)\\
\textbf{Adaptive Capacity}: 2 PCs retained (69\% variance explained)
\end{quote}
Higher variance explained means the pillar's main pattern does a good job
of summarizing all the included variables. Very low variance explained
(say, below 50\%) is a hint that your included variables are pulling in
too many different directions and the summary score may be misleading.

%% ==================================================================
\section{Example Workflows}
%% ==================================================================

\subsection{Workflow 1: ``Where should we put the next cooling center?''}

\begin{enumerate}[leftmargin=*]
  \item Start with default variable selections.
  \item Under \textcolor{adaptive}{Adaptive Capacity}, \textit{uncheck}
        ``AC Prevalence'' and ``Prob.\ Central AC.'' You want a view that
        penalizes the lack of AC rather than treating homes with AC
        equivalently.
  \item Click \textbf{Compute PCA}.
  \item Select the \textbf{Composite Risk} layer.
  \item Zoom to the tracts coloring darkest red.
  \item Click a few to confirm the ``dominant driver'' is
        adaptive-capacity-related. If it's sensitivity or exposure instead,
        a cooling center alone won't fully address the problem.
\end{enumerate}

\subsection{Workflow 2: ``Which neighborhoods will warm the most by mid-century?''}

\begin{enumerate}[leftmargin=*]
  \item Click \textbf{SSP3-7.0} (or 5-8.5 for a worst case).
  \item Keep only the \textcolor{exposure}{Exposure} variables checked.
  \item Click \textbf{Compute PCA}.
  \item Compare side-by-side with a second browser window showing
        \textbf{Present Day} to see where warming is concentrated.
\end{enumerate}

\subsection{Workflow 3: ``Make the case for tree planting in one tract.''}

\begin{enumerate}[leftmargin=*]
  \item Use default selections and compute PCA.
  \item Click your tract of interest.
  \item Note its composite score, its dominant driver, and how its tree
        canopy percentile compares to neighboring tracts.
  \item Screenshot and share. The methodology in Sections~\ref{sec:pillars}
        and~\ref{sec:pca} of this guide is your documentation.
\end{enumerate}

%% ==================================================================
\section{Interpretation and Caveats}
%% ==================================================================

\subsection{What the Scores Mean --- and Don't Mean}

\begin{itemize}[leftmargin=*]
  \item The scores are \textbf{relative}, not absolute. A ``75'' means
        the tract is in the 75th percentile among San Diego County tracts,
        not that 75\% of residents will be harmed.
  \item The scores are \textbf{descriptive}, not predictive. They
        summarize current vulnerability patterns; they do not forecast
        specific heat-related illnesses or deaths.
  \item The scores depend on \textbf{your variable choices}. Different
        choices give different scores. There is no single ``correct''
        answer --- that is the point of a DIY tool.
\end{itemize}

\subsection{Things That Will Break the Analysis}

\begin{itemize}[leftmargin=*]
  \item \textbf{Too few variables.} A pillar with one or zero checked
        variables can't produce a meaningful PCA. The tool will show
        ``no variables'' in the results summary.
  \item \textbf{Highly redundant variables.} Including 10 near-identical
        temperature measures won't add information; PCA will collapse them
        into one component, and the pillar will be dominated by that one
        theme.
  \item \textbf{Mixing directions without thinking.} If you move
        ``Life Expectancy'' (protective) into Sensitivity, make sure you
        know what you're doing. The defaults already handle sign
        direction for you; moving variables between pillars requires care.
\end{itemize}

\subsection{Data Vintage and Coverage}

\begin{itemize}[leftmargin=*]
  \item Temperature data: weather station observations (MesoWest/Synoptic,
        CIMIS) from recent years.
  \item Future projections: LOCA2 downscaled CMIP6 ensemble mean,
        2035--2064 window.
  \item Health data: CDC PLACES (most recent release), SVI 2022.
  \item Socioeconomic data: Healthy Places Index (most recent release),
        American Community Survey 5-year estimates.
  \item AC prevalence: Romitti et al.\ and Ahn et al.\ modeled estimates.
\end{itemize}

Every variable's row has a \texttt{src} link to the original dataset ---
follow it before citing a specific number.

\subsection{What This Tool Is Not}

\begin{itemize}[leftmargin=*]
  \item \textbf{Not a real-time forecast.} For current heat warnings,
        use the National Weather Service.
  \item \textbf{Not a substitute for local knowledge.} A tract-level
        score can miss the block with the mobile home park next to the
        freeway. Ground-truthing matters.
  \item \textbf{Not a replacement for engagement.} The communities
        shown as highest risk are the ones whose input should shape
        any intervention.
\end{itemize}

%% ==================================================================
\section{Frequently Asked Questions}
%% ==================================================================

\paragraph{Can I use this outside San Diego County?}
Not with the current data. The pipeline itself is portable --- every
variable comes from a freely available, nationally collected dataset ---
but someone would need to re-run the data assembly for a new county.

\paragraph{Can I save or export my map?}
The current version of the tool runs entirely in your browser. Take
screenshots of the map and the results summary to document a specific
configuration. A formal export feature is on the roadmap.

\paragraph{Why are some tracts blank?}
A tract is shown blank when one or more of the underlying variables you
selected is missing for that tract. Try removing a variable with spotty
coverage to see if the blanks fill in.

\paragraph{What's the difference between ``Composite Risk'' and a single
pillar's PC1?}
Composite Risk combines all three pillars. A pillar's PC1 shows only the
dominant pattern within that one pillar --- useful when you want to see,
for example, which neighborhoods are hottest regardless of who lives in
them.

\paragraph{Who do I contact for help or to report an issue?}
See the project's \texttt{README.md} for current maintainer contact
information.

%% ==================================================================
\section{Glossary}
%% ==================================================================

\begin{description}[style=nextline]
  \item[Adaptive Capacity] Resources and conditions that help a community
        cope with heat --- AC, tree canopy, income, healthcare access.
  \item[Census Tract] A small, relatively permanent geographic subdivision
        used by the U.S.\ Census Bureau. There are roughly 630 tracts in
        San Diego County, each averaging 4,000 residents.
  \item[Composite Risk] The tool's final 0--100 score combining all three
        pillars.
  \item[Dominant Driver] The pillar contributing most to a given tract's
        composite risk.
  \item[Eigenvalue] A number PCA produces that describes how important a
        given pattern is. Patterns with eigenvalues above 1 are kept; the
        rest are dropped.
  \item[Exposure] The physical heat experienced in a place ---
        temperature, humidity.
  \item[Kaiser Criterion] The rule used to keep only PCA components with
        eigenvalues greater than 1.
  \item[LOCA2] A high-resolution downscaled climate projection dataset,
        used here for future temperature changes.
  \item[PCA (Principal Component Analysis)] A statistical method for
        finding the most important hidden patterns in a set of correlated
        variables.
  \item[Percentile Rank] A score from 0 to 100 representing where a tract
        sits relative to all other tracts in the county.
  \item[Pillar] One of the three branches of the risk framework ---
        Exposure, Sensitivity, or Adaptive Capacity.
  \item[Sensitivity] Population characteristics that make people more
        likely to be harmed by heat --- age, chronic disease, social
        vulnerability.
  \item[SSP] ``Shared Socioeconomic Pathway'' --- a scenario for future
        emissions used by climate models. Higher numbers mean more warming.
  \item[SVI (Social Vulnerability Index)] The CDC's index of community
        vulnerability to disasters, based on Census/ACS variables.
  \item[Variance Explained] The percentage of information in the original
        variables captured by a given PCA component. Higher is better.
\end{description}

\vspace{1em}
\hrule
\vspace{1em}

\noindent\textit{This guide describes the San Diego Heat Risk Explorer
(file: \texttt{figures/main/diy\_pca.html}) produced by
\texttt{scripts/build\_diy\_pca\_viewer.py}. For the technical methodology
behind the scores, see \texttt{docs/risk\_index/methods.tex}.}

\end{document}
